% 第1章《引论》习题解答
% 对应 chapters/ch01.tex

\chapter*{第1章：引论习题解答}
\addcontentsline{toc}{chapter}{第1章：引论习题解答}
\section*{1.3.2\quad 练习题}

\begin{xexercise}[1]
关于 Bernoulli 不等式的推广。
\end{xexercise}
\begin{xsolution}
\begin{enumerate}[label=(\arabic*)]
  \item 当 $-2\le h\le-1$ 时，$-1\le1+h\le0$。对任意 $n\in\NN_+$，
  \[
    (1+h)^n\ge-\abs{1+h}=1+h.
  \]
  又因 $h\le0$ 且 $n\ge1$，有 $h\ge nh$，故
  \[
    (1+h)^n\ge1+h\ge1+nh.
  \]
  因而 Bernoulli 不等式在该区间仍成立。

  \item 当 $h\ge0$ 时，由二项式展开
  \[
    (1+h)^n
    =\sum_{j=0}^n\binom{n}{j}h^j
    =1+nh+\frac{n(n-1)}2h^2+\cdots+h^n.
  \]
  各项均非负，所以
  \[
    (1+h)^n\ge\frac{n(n-1)}2h^2.
  \]
  同一理由给出更一般的结论：对任意整数 $m$，$0\le m\le n$，均有
  \[
    (1+h)^n\ge\sum_{j=0}^m\binom{n}{j}h^j
    \ge\binom{n}{m}h^m,
    \qquad h\ge0.
  \]

  \item 对 $n$ 作数学归纳。$n=1$ 时为等式。设结论对 $n=k$ 成立，即
  \[
    \prod_{i=1}^k(1+a_i)\ge1+\sum_{i=1}^k a_i.
  \]
  因 $a_{k+1}>-1$，有 $1+a_{k+1}>0$，故
  \[
  \begin{aligned}
    \prod_{i=1}^{k+1}(1+a_i)
    &\ge\left(1+\sum_{i=1}^k a_i\right)(1+a_{k+1})\\
    &=1+\sum_{i=1}^{k+1}a_i
      +a_{k+1}\sum_{i=1}^k a_i.
  \end{aligned}
  \]
  题设各 $a_i$ 同号，所以
  \[
    a_{k+1}\sum_{i=1}^k a_i\ge0.
  \]
  因而
  \[
    \prod_{i=1}^{k+1}(1+a_i)\ge1+\sum_{i=1}^{k+1}a_i,
  \]
  归纳完成。
\end{enumerate}
\end{xsolution}

\begin{xexercise}[2]
关于阶乘 $n!$ 的几个不等式。
\end{xexercise}
\begin{xsolution}
设 $n>1$。
\begin{enumerate}[label=(\arabic*)]
  \item 对正数 $1,2,\ldots,n$ 使用算术平均值—几何平均值不等式，
  \[
    (n!)^{1/n}
    <\frac{1+2+\cdots+n}{n}
    =\frac{n+1}{2}.
  \]
  这里严格不等号成立，是因为 $n>1$ 时这些数不全相等。两边取 $n$ 次幂即得
  \[
    n!<\left(\frac{n+1}{2}\right)^n.
  \]

  \item 由
  \[
    (n!)^2=\prod_{k=1}^n k(n+1-k)
  \]
  对 $n$ 个正数 $k(n+1-k)$ 使用算术平均值—几何平均值不等式，得到
  \[
    (n!)^{2/n}
    \le\frac1n\sum_{k=1}^n k(n+1-k).
  \]
  计算右端：
  \[
  \begin{aligned}
    \frac1n\sum_{k=1}^n k(n+1-k)
    &=\frac1n\left((n+1)\sum_{k=1}^n k-\sum_{k=1}^n k^2\right)\\
    &=\frac{(n+1)(n+2)}6.
  \end{aligned}
  \]
  因此实际上有更强的估计
  \[
    n!\le\left(\sqrt{\frac{(n+1)(n+2)}6}\right)^n.
  \]
  又因为 $n+1<n+2$，故
  \[
    \sqrt{\frac{(n+1)(n+2)}6}<\frac{n+2}{\sqrt6},
  \]
  从而
  \[
    n!<\left(\frac{n+2}{\sqrt6}\right)^n.
  \]

  \item 比较题中两上界的底数
  \[
    A_n=\frac{n+1}{2},\qquad B_n=\frac{n+2}{\sqrt6}.
  \]
  因二者均为正数，$B_n<A_n$ 等价于
  \[
    2(n+2)<\sqrt6(n+1),
  \]
  进一步等价于
  \[
    n^2-2n-5>0.
  \]
  对整数 $n>1$，这恰在 $n\ge4$ 时成立。因此：$n=2,3$ 时 (1) 的上界较小，故 (1) 较优；$n\ge4$ 时 (2) 的上界较小，故 (2) 较优。

  还可注意到，(2) 的证明中得到的中间估计
  \[
    n!\le\left(\sqrt{\frac{(n+1)(n+2)}6}\right)^n
  \]
  对每个 $n>1$ 都比 (1) 更强，因为
  \[
    \sqrt{\frac{(n+1)(n+2)}6}<\frac{n+1}{2}
  \]
  等价于 $n>1$。

  \item 对任意实数 $r$，各数 $1^r,2^r,\ldots,n^r$ 都为正数。由算术平均值—几何平均值不等式，
  \[
    \left(\prod_{k=1}^n k^r\right)^{1/n}
    \le\frac1n\sum_{k=1}^n k^r.
  \]
  即
  \[
    (n!)^{r/n}\le\frac1n\sum_{k=1}^n k^r.
  \]
  两边取 $n$ 次幂，得到
  \[
    (n!)^r\le\frac1{n^n}\left(\sum_{k=1}^n k^r\right)^n.
  \]
\end{enumerate}
\end{xsolution}

\begin{xexercise}[3]
证明几何平均值—调和平均值不等式。
\end{xexercise}
\begin{xsolution}
由 $a_k>0$，可对正数 $1/a_1,\ldots,1/a_n$ 使用算术平均值—几何平均值不等式：
\[
  \frac1n\sum_{k=1}^n\frac1{a_k}
  \ge
  \left(\prod_{k=1}^n\frac1{a_k}\right)^{1/n}
  =\left(\prod_{k=1}^n a_k\right)^{-1/n}.
\]
两边均为正数，取倒数后不等号反向，故
\[
  \frac{n}{\displaystyle\sum_{k=1}^n\frac1{a_k}}
  \le
  \left(\prod_{k=1}^n a_k\right)^{1/n}.
\]
这就是所要证明的不等式。等号成立当且仅当 $1/a_1=\cdots=1/a_n$，亦即 $a_1=\cdots=a_n$。
\end{xsolution}

\begin{xexercise}[4]
证明三个非负数的几何平均值、平方平均式与算术平均值之间的不等式。
\end{xexercise}
\begin{xsolution}
对非负数 $ab,bc,ca$ 使用算术平均值—几何平均值不等式，
\[
  \frac{ab+bc+ca}{3}
  \ge\sqrt[3]{(ab)(bc)(ca)}
  =(abc)^{2/3}.
\]
两边开平方，得到
\[
  \sqrt{\frac{ab+bc+ca}{3}}\ge\sqrt[3]{abc}.
\]
另一方面，
\[
  (a+b+c)^2-3(ab+bc+ca)
  =\frac12\bigl((a-b)^2+(b-c)^2+(c-a)^2\bigr)\ge0.
\]
于是
\[
  \frac{ab+bc+ca}{3}\le\frac{(a+b+c)^2}{9}.
\]
因两边非负，开平方即得
\[
  \sqrt{\frac{ab+bc+ca}{3}}\le\frac{a+b+c}{3}.
\]
综上，
\[
  \sqrt[3]{abc}\le\sqrt{\frac{ab+bc+ca}{3}}\le\frac{a+b+c}{3}.
\]
两处同时取等号当且仅当 $a=b=c$。
\end{xsolution}

\begin{xexercise}[5]
证明题中四组不等式，并回答第 (2) 小问。
\end{xexercise}
\begin{xsolution}
\begin{enumerate}[label=(\arabic*)]
  \item 由三角形不等式
  \[
    \abs{a}=\abs{(a-b)+b}\le\abs{a-b}+\abs{b},
  \]
  得
  \[
    \abs{a-b}\ge\abs{a}-\abs{b}.
  \]
  交换 $a,b$，又有
  \[
    \abs{a-b}=\abs{b-a}\ge\abs{b}-\abs{a}.
  \]
  因而 $\abs{a-b}$ 同时不小于 $\abs{a}-\abs{b}$ 与其相反数，所以
  \[
    \abs{a-b}\ge\abs{\abs{a}-\abs{b}}.
  \]

  \item 反复使用三角形不等式，
  \[
    \abs{\sum_{k=1}^n a_k}\le\sum_{k=1}^n\abs{a_k}.
  \]
  又由
  \[
    a_1=\sum_{k=1}^n a_k-\sum_{k=2}^n a_k
  \]
  得
  \[
    \abs{a_1}
    \le\abs{\sum_{k=1}^n a_k}+\sum_{k=2}^n\abs{a_k},
  \]
  即
  \[
    \abs{a_1}-\sum_{k=2}^n\abs{a_k}
    \le\abs{\sum_{k=1}^n a_k}.
  \]
  因此题中双边不等式成立。

  左边一般不能改为
  \[
    \abs{\abs{a_1}-\sum_{k=2}^n\abs{a_k}}.
  \]
  例如取 $n=3$，$a_1=1$，$a_2=10$，$a_3=-10$，则
  \[
    \abs{\abs{a_1}-\abs{a_2}-\abs{a_3}}=19,
    \qquad
    \abs{a_1+a_2+a_3}=1,
  \]
  所以这种加强不成立。

  \item 记 $x=\abs{a}$，$y=\abs{b}$。由三角形不等式及函数 $t/(1+t)$ 在 $t\ge0$ 上单调增加，
  \[
    \frac{\abs{a+b}}{1+\abs{a+b}}
    \le\frac{x+y}{1+x+y}.
  \]
  另一方面，
  \[
  \begin{aligned}
    \frac{x}{1+x}+\frac{y}{1+y}-\frac{x+y}{1+x+y}
    &=\frac{xy(2+x+y)}{(1+x)(1+y)(1+x+y)}\\
    &\ge0.
  \end{aligned}
  \]
  因此
  \[
    \frac{\abs{a+b}}{1+\abs{a+b}}
    \le\frac{\abs{a}}{1+\abs{a}}+\frac{\abs{b}}{1+\abs{b}}.
  \]

  \item 由二项式展开与三角形不等式，
  \[
  \begin{aligned}
    \abs{(a+b)^n-a^n}
    &=\abs{\sum_{k=1}^n\binom{n}{k}a^{n-k}b^k}\\
    &\le\sum_{k=1}^n\binom{n}{k}\abs{a}^{n-k}\abs{b}^k\\
    &=(\abs{a}+\abs{b})^n-\abs{a}^n.
  \end{aligned}
  \]
\end{enumerate}
\end{xsolution}

\begin{xexercise}[6]
按四个提示给出 Cauchy 不等式的不同证明。
\end{xexercise}
\begin{xsolution}
记
\[
  S_n=\sum_{k=1}^n a_kb_k,
  \qquad
  A_n=\sum_{k=1}^n a_k^2,
  \qquad
  B_n=\sum_{k=1}^n b_k^2.
\]
要证的是 $\abs{S_n}^2\le A_nB_n$。

\begin{enumerate}[label=(\arabic*)]
  \item 用数学归纳法。$n=1$ 时为等式。设结论对 $n$ 成立。令
  \[
    A=\sqrt{A_n},\qquad B=\sqrt{B_n},\qquad a=a_{n+1},\qquad b=b_{n+1}.
  \]
  由归纳假设及三角形不等式，
  \[
    \abs{S_{n+1}}\le AB+\abs{ab}.
  \]
  而
  \[
  \begin{aligned}
    &(A^2+a^2)(B^2+b^2)-(AB+\abs{ab})^2\\
    &\qquad=A^2b^2+B^2a^2-2AB\abs{ab}
    =(A\abs{b}-B\abs{a})^2\ge0.
  \end{aligned}
  \]
  所以
  \[
    AB+\abs{ab}\le\sqrt{A^2+a^2}\sqrt{B^2+b^2}.
  \]
  因而 Cauchy 不等式对 $n+1$ 成立，归纳完成。

  \item 由题给 Lagrange 恒等式，
  \[
    A_nB_n-\left(\sum_{k=1}^n\abs{a_kb_k}\right)^2
    =\frac12\sum_{k=1}^n\sum_{i=1}^n
    (\abs{a_k}\abs{b_i}-\abs{a_i}\abs{b_k})^2\ge0.
  \]
  故
  \[
    \sum_{k=1}^n\abs{a_kb_k}\le\sqrt{A_nB_n}.
  \]
  再由
  \[
    \abs{S_n}\le\sum_{k=1}^n\abs{a_kb_k},
  \]
  即得 Cauchy 不等式。

  \item 若 $A_n=0$ 或 $B_n=0$，结论立即成立。以下设 $A_nB_n>0$，令
  \[
    \alpha_k=\frac{a_k}{\sqrt{A_n}},
    \qquad
    \beta_k=\frac{b_k}{\sqrt{B_n}}.
  \]
  则 $\sum\alpha_k^2=\sum\beta_k^2=1$。由
  \[
    \abs{AB}\le\frac{A^2+B^2}{2}
  \]
  得
  \[
  \begin{aligned}
    \frac{\abs{S_n}}{\sqrt{A_nB_n}}
    &\le\sum_{k=1}^n\abs{\alpha_k\beta_k}\\
    &\le\frac12\sum_{k=1}^n(\alpha_k^2+\beta_k^2)=1.
  \end{aligned}
  \]
  故 $\abs{S_n}\le\sqrt{A_nB_n}$。

  \item 令
  \[
    c_k=a_k^2-b_k^2+2\ii\abs{a_kb_k}.
  \]
  则
  \[
    \abs{c_k}
    =\sqrt{(a_k^2-b_k^2)^2+4a_k^2b_k^2}
    =a_k^2+b_k^2.
  \]
  由题给三角形不等式，
  \[
    \abs{\sum_{k=1}^n c_k}\le A_n+B_n.
  \]
  左边平方为
  \[
    (A_n-B_n)^2+4\left(\sum_{k=1}^n\abs{a_kb_k}\right)^2.
  \]
  因此
  \[
    (A_n-B_n)^2+4\left(\sum_{k=1}^n\abs{a_kb_k}\right)^2
    \le(A_n+B_n)^2,
  \]
  即
  \[
    \left(\sum_{k=1}^n\abs{a_kb_k}\right)^2\le A_nB_n.
  \]
  最后用 $\abs{S_n}\le\sum\abs{a_kb_k}$，即得所需结论。
\end{enumerate}
\end{xsolution}

\begin{xexercise}[7]
用向前—向后数学归纳法证明 Ky Fan 不等式。
\end{xexercise}
\begin{xsolution}
对 $0<x_i\le\dfrac12$，记
\[
  A_n=\frac1n\sum_{i=1}^n x_i,
  \qquad
  G_n=\left(\prod_{i=1}^n x_i\right)^{1/n},
\]
并记
\[
  A_n'=1-A_n=\frac1n\sum_{i=1}^n(1-x_i),
  \qquad
  G_n'=\left(\prod_{i=1}^n(1-x_i)\right)^{1/n}.
\]
题中不等式等价于
\[
  \frac{G_n}{A_n}\le\frac{G_n'}{A_n'}.
\]
记此命题为 $P_n$。

先证 $P_2$。令 $s=x_1+x_2$，$p=x_1x_2$。因 $0<x_i\le1/2$，有 $s\le1$，且
\[
  s^2-4p=(x_1-x_2)^2\ge0.
\]
$P_2$ 平方并清去正分母后等价于
\[
  p(2-s)^2\le(1-s+p)s^2.
\]
而两边之差为
\[
  (1-s+p)s^2-p(2-s)^2
  =(s^2-4p)(1-s)\ge0,
\]
故 $P_2$ 成立。

下面证明“向前”步骤：若 $P_n$ 成立，则 $P_{2n}$ 成立。把 $2n$ 个数分成两组，每组 $n$ 个。两组的算术平均值、几何平均值分别记为
\[
  A_1,G_1,A_1',G_1',
  \qquad
  A_2,G_2,A_2',G_2'.
\]
由 $P_n$，
\[
  \frac{G_j}{G_j'}\le\frac{A_j}{A_j'},\qquad j=1,2.
\]
于是
\[
  \frac{G_{2n}}{G_{2n}'}
  =\sqrt{\frac{G_1G_2}{G_1'G_2'}}
  \le\sqrt{\frac{A_1A_2}{A_1'A_2'}}.
\]
又因 $0<A_1,A_2\le1/2$，可对两个数 $A_1,A_2$ 使用已经证明的 $P_2$，得到
\[
  \sqrt{\frac{A_1A_2}{A_1'A_2'}}
  \le\frac{(A_1+A_2)/2}{(A_1'+A_2')/2}
  =\frac{A_{2n}}{A_{2n}'}.
\]
故 $P_{2n}$ 成立。从 $P_2$ 出发，依次得到 $P_4,P_8,\ldots,P_{2^m}$。

再证明“向后”步骤：若 $P_n$ 对某个 $n>1$ 成立，则 $P_{n-1}$ 成立。给定 $n-1$ 个数 $x_1,\ldots,x_{n-1}$，令
\[
  x_n=A_{n-1}=\frac1{n-1}\sum_{i=1}^{n-1}x_i.
\]
仍有 $0<x_n\le1/2$，且加入这个数后
\[
  A_n=A_{n-1},\qquad A_n'=A_{n-1}'.
\]
由 $P_n$，
\[
  \frac{(G_{n-1}^{n-1}A_{n-1})^{1/n}}{A_{n-1}}
  \le
  \frac{((G_{n-1}')^{n-1}A_{n-1}')^{1/n}}{A_{n-1}'}.
\]
两边取 $n$ 次幂并约去正因子，得到
\[
  \left(\frac{G_{n-1}}{A_{n-1}}\right)^{n-1}
  \le
  \left(\frac{G_{n-1}'}{A_{n-1}'}\right)^{n-1},
\]
即 $P_{n-1}$ 成立。

任取正整数 $n$，选择 $m$ 使 $2^m\ge n$。由向前步骤知 $P_{2^m}$ 成立，再反复使用向后步骤即可得到 $P_n$。因此题中不等式对所有 $n$ 成立。
\end{xsolution}

\begin{xexercise}[8]
设 $a,c,g,t\ge0$，$a+c+g+t=1$，证明平方和的下界并确定等号条件。
\end{xexercise}
\begin{xsolution}
由 Cauchy 不等式，
\[
  (a+c+g+t)^2
  \le(1^2+1^2+1^2+1^2)(a^2+c^2+g^2+t^2).
\]
利用 $a+c+g+t=1$，得到
\[
  1\le4(a^2+c^2+g^2+t^2),
\]
即
\[
  a^2+c^2+g^2+t^2\ge\frac14.
\]
Cauchy 不等式取等号当且仅当 $(a,c,g,t)$ 与 $(1,1,1,1)$ 成比例。再结合四数之和为 $1$，可知等号成立当且仅当
\[
  a=c=g=t=\frac14.
\]
\end{xsolution}

\section*{1.4\quad 逻辑符号与对偶法则：练习题}

\begin{xexercise}[1.4 练习题]
以正面方式写出原章所列八个命题或叙述的否定。
\end{xexercise}
\begin{xsolution}
以下均按原命题的量词结构逐层取否定。
\begin{enumerate}[label=(\arabic*)]
  \item “数集 $A$ 有上界”可写成
  \[
    \exists M\in\RR,\ \forall x\in A,\quad x\le M.
  \]
  因此其否定为
  \[
    \forall M\in\RR,\ \exists x\in A,\quad x>M.
  \]

  \item “数集 $A$ 的最小值是 $b$”可写成
  \[
    b\in A,
    \qquad
    \forall x\in A,\quad b\le x.
  \]
  因此其否定为
  \[
    b\notin A
    \quad\text{或}\quad
    \exists x\in A,\quad x<b.
  \]

  \item “$f$ 是区间 $(a,b)$ 上的单调增加函数”可写为
  \[
    \forall x_1,x_2\in(a,b),\quad
    x_1<x_2\Longrightarrow f(x_1)\le f(x_2).
  \]
  因此其否定为
  \[
    \exists x_1,x_2\in(a,b),\quad
    x_1<x_2,\qquad f(x_1)>f(x_2).
  \]

  \item “$f$ 是区间 $(a,b)$ 上的单调函数”表示 $f$ 单调增加或单调减少。其否定要求两种情形都不成立，即
  \[
  \begin{gathered}
    \exists x_1,x_2\in(a,b),\quad x_1<x_2,\quad f(x_1)>f(x_2),\\
    \exists y_1,y_2\in(a,b),\quad y_1<y_2,\quad f(y_1)<f(y_2).
  \end{gathered}
  \]
  两对点不必相同。

  \item $A\subset B$ 表示
  \[
    \forall x,\quad x\in A\Longrightarrow x\in B.
  \]
  因此其否定为
  \[
    \exists x,\quad x\in A,\qquad x\notin B.
  \]

  \item $A-B\ne\varnothing$ 等价于
  \[
    \exists x,\quad x\in A,\qquad x\notin B.
  \]
  因此其否定为
  \[
    \forall x\in A,\quad x\in B,
  \]
  即 $A-B=\varnothing$。

  \item “数列 $\{x_n\}$ 是无穷小量”即
  \[
    \forall\varepsilon>0,\ \exists N,\ \forall n>N,
    \quad \abs{x_n}<\varepsilon.
  \]
  其否定为
  \[
    \exists\varepsilon_0>0,\ \forall N,\ \exists n>N,
    \quad \abs{x_n}\ge\varepsilon_0.
  \]

  \item “数列 $\{x_n\}$ 是正无穷大量”即
  \[
    \forall M>0,\ \exists N,\ \forall n>N,
    \quad x_n>M.
  \]
  其否定为
  \[
    \exists M_0>0,\ \forall N,\ \exists n>N,
    \quad x_n\le M_0.
  \]
\end{enumerate}
\end{xsolution}
